\documentclass[10pt,twocolumn]{article}

% ── Packages ──────────────────────────────────────────────────────────────────
\usepackage[a4paper, top=2cm, bottom=2cm, left=1.8cm, right=1.8cm]{geometry}
\usepackage{times}
\usepackage{graphicx}
\usepackage{booktabs}
\usepackage{tabularx}
\usepackage{amsmath}
\usepackage{hyperref}
\usepackage{enumitem}
\usepackage{caption}
\usepackage{subcaption}
\usepackage{titlesec}
\usepackage{abstract}
\usepackage{fancyhdr}
\usepackage{microtype}
\usepackage{stfloats}
\usepackage[T1]{fontenc}
\usepackage[utf8]{inputenc}

% ── Section formatting ────────────────────────────────────────────────────────
\titleformat{\section}{\normalfont\bfseries\normalsize}{\thesection.}{0.5em}{}
\titleformat{\subsection}{\normalfont\bfseries\small}{\thesubsection.}{0.5em}{}
\titlespacing{\section}{0pt}{6pt}{3pt}
\titlespacing{\subsection}{0pt}{4pt}{2pt}

% ── Abstract formatting ───────────────────────────────────────────────────────
\renewcommand{\abstractnamefont}{\normalfont\bfseries\small}
\renewcommand{\abstracttextfont}{\normalfont\small}
\setlength{\absleftindent}{0pt}
\setlength{\absrightindent}{0pt}

% ── Header / Footer ───────────────────────────────────────────────────────────
\pagestyle{fancy}
\fancyhf{}
\rhead{\small Cloud Computing --- Sapienza University of Rome, 2026}
\lhead{\small Option 2: Image Processing on AWS Lambda}
\cfoot{\small\thepage}
\renewcommand{\headrulewidth}{0.4pt}

% ── Caption formatting ────────────────────────────────────────────────────────
\captionsetup{font=small, labelfont=bf, skip=4pt}
\captionsetup[subfigure]{font=small, labelfont=bf, skip=3pt}

% ── Hyperlinks ────────────────────────────────────────────────────────────────
\hypersetup{colorlinks=true, linkcolor=blue, urlcolor=blue, citecolor=blue}

% ── Misc ──────────────────────────────────────────────────────────────────────
\setlength{\columnsep}{0.6cm}
\setlist[itemize]{noitemsep, topsep=2pt, leftmargin=*}
\setlist[enumerate]{noitemsep, topsep=2pt, leftmargin=*}

% ══════════════════════════════════════════════════════════════════════════════
\begin{document}

% ── Title block ───────────────────────────────────────────────────────────────
\twocolumn[{%
\begin{center}
  {\Large\bfseries Image Processing on AWS Lambda: \\[4pt]
   Performance Evaluation and Cost Analysis}\\[8pt]
  {\normalsize
    Nicol\'as Pacheco \quad Prajwal Shinde \quad Mathias Tirado\\[2pt]
    \textit{Cloud Computing, Prof.\ Emiliano Casalicchio}\\
    \textit{Sapienza Universit\`a di Roma --- June 2026}
  }
\end{center}
\vspace{6pt}
\begin{onecolabstract}
We deploy three image-processing operations (\textit{resize}, \textit{grayscale
conversion}, and \textit{edge detection}) as independent AWS Lambda functions
exposed through Amazon API Gateway, and evaluate their performance, scalability,
and operational cost. A factorial load-testing matrix of 3~operations
$\times$ 3~image sizes $\times$ 4~concurrency levels $\times$ 2~repetitions
(72~scenarios) is executed with Locust from a team member's personal laptop reaching the
Lambda endpoints over the public Internet, while five CloudWatch metrics per
Lambda capture the server-side picture. Lambda processes
every request in under 1.5\,s with zero errors and zero throttles across
52{,}272~invocations, yet end users observe up to 46\,s p95 response time and
27\,\% HTTP failures at 100~concurrent users. All failures are rejected
upstream by API Gateway before reaching Lambda. A six-month cost projection places the
break-even point against an always-on EC2~t3.small at roughly 7\,M
invocations per month.
\end{onecolabstract}
\vspace{10pt}
}]

% ══════════════════════════════════════════════════════════════════════════════
\section{Introduction}
\label{sec:intro}

Serverless computing platforms such as AWS Lambda are increasingly adopted for
event-driven, stateless workloads because they eliminate infrastructure
management and bill only for actual execution time~\cite{aws:lambda-pricing}.
Image processing---characterised by variable payload sizes, purely stateless
computation, and bursty request patterns---is a natural fit for this model, and
serves as a concrete, measurable workload for evaluating serverless scalability.

This report describes the design, deployment, and empirical evaluation of an
image-processing pipeline built on AWS Lambda following Option~2 of the course
project. Three operations---resize, grayscale conversion, and edge
detection---each run as a separate Lambda function exposed through Amazon API
Gateway. The central evaluation question is \textit{how a serverless deployment
scales from one to one hundred concurrent users}, and crucially, how server-side
Lambda metrics relate to the latency that end users actually perceive.

We executed 72 Locust load scenarios, cross-referenced client-side measurements
with five per-Lambda CloudWatch metrics (Section~\ref{sec:server}--\ref{sec:cross}),
and projected the six-month operational cost against an EC2 baseline
(Section~\ref{sec:cost}). Our central finding: Lambda reported perfect server-side health while end
users experienced up to 46\,s response times and 27\,\% HTTP failures. This
illustrates why monitoring a single layer of a multi-tier architecture
produces a misleading picture of system health.

% ══════════════════════════════════════════════════════════════════════════════
\section{System Architecture}
\label{sec:arch}

The system follows a \textbf{microservices} pattern: each operation runs in its
own Lambda function with its own API Gateway endpoint. All three Lambdas share
the same runtime---Python~3.11, x86\_64, 512\,MB memory, 30-second function
timeout---and use the public Klayers Pillow~11 layer for image
manipulation~\cite{pillow:docs}. API Gateway HTTP~API endpoints retain the
default 29-second integration timeout, which becomes the binding constraint at
high concurrency (Section~\ref{sec:cross}).

\subsection{Per-operation algorithms}

Each operation is a thin wrapper around Pillow primitives, intentionally keeping
application logic minimal so that measurement focuses on infrastructure
behaviour rather than algorithm tuning:

\begin{itemize}
  \item \textbf{Resize} --- LANCZOS resampling scales the input to a target
        width of 400\,px while preserving aspect ratio. LANCZOS is the
        highest-quality interpolation filter in Pillow and is chosen to
        maximise per-invocation CPU work, making duration differences
        across image sizes more pronounced.
  \item \textbf{Grayscale} --- \texttt{convert("L")} converts to single-channel
        8-bit luminance, then re-encodes as JPEG at quality~85. The lightest
        of the three operations: each output pixel requires only a weighted sum
        of the three input colour channels.
  \item \textbf{Edge detection} --- a 3$\times$3 SMOOTH averaging filter is
        first applied to reduce noise, followed by the FIND\_EDGES convolution
        kernel on the grayscale image~\cite{pillow:docs}. Two sequential
        convolution passes make this the most CPU-intensive operation.
\end{itemize}

\subsection{Microservices vs.\ monolithic Lambda}

We initially prototyped a single monolithic Lambda in which the operation was
selected by a request parameter, before migrating to the three-function
decomposition. A microservice is characterised by \textit{single
responsibility}, \textit{independent deployability}, \textit{independent
scalability} and \textit{isolated failure modes}---properties that map
naturally to AWS Lambda, where each function is by construction a separately
deployed and separately scaled execution unit with its own IAM role, code
package, and CloudWatch namespace. Our three-function decomposition therefore
aligns Lambda's runtime model with the microservices architectural style; a
monolithic Lambda that routes between unrelated capabilities would borrow
Lambda's deployment model but lose those architectural properties.

The decisive consideration for this project, however, is \textit{per-function
observability}: CloudWatch publishes metrics in a per-function namespace, so a
monolithic Lambda aggregates Duration and Concurrent Executions across all
three operations, making per-operation attribution impossible. With independent
functions we can report that grayscale achieves p95\,=\,432\,ms while edge
detection reaches p95\,=\,1{,}054\,ms (Table~\ref{tab:cloudwatch}). This 2.4$\times$
gap would be unextractable from a shared namespace. The trade-off is a
tripled deployment footprint (three functions, three endpoints, three IAM
policies), which is acceptable given the measurement benefit.

% ══════════════════════════════════════════════════════════════════════════════
\section{Experimental Design}
\label{sec:design}

\subsection{Factorial matrix}

The test matrix crosses three orthogonal factors:
\textbf{operation} (resize, grayscale, edge) $\times$
\textbf{image size} (small $\approx$55\,KB, medium $\approx$400\,KB,
large $\approx$1.3\,MB) $\times$
\textbf{concurrent users} (1, 10, 50, 100),
each combination repeated twice for variance estimation.
Total: $3 \times 3 \times 4 \times 2 = 72$ scenarios, each running for
2\,minutes at steady state after a 10-second ramp-up. The original plan included
200 users and 5-minute runs; both were reduced to stay within the \$50 Vocareum
credit budget. The 2-minute duration was sufficient for every scenario to reach
a stable operating point, as confirmed by flat user-count curves in the
per-scenario history files.

Test images were generated programmatically with Pillow using random-noise pixel
fills. Solid-colour images were explicitly avoided because JPEG compression
collapses them to a few kilobytes regardless of canvas dimensions, eliminating
the payload-size variable. Random noise represents worst-case JPEG entropy and
produces file sizes that reliably match the three target size categories.

\subsection{Data collection}

Two complementary data sources were used in parallel throughout every scenario:

\begin{itemize}
  \item \textbf{Locust~\cite{locust:docs} (client-side)} --- per-scenario
        \texttt{stats.csv} with aggregated p50/p95/p99 response times,
        throughput, and failure counts, plus a per-second
        \texttt{stats\_history.csv} timeline. Captures the full end-to-end
        latency including network transit and all queuing delays between the
        load generator and Lambda.
  \item \textbf{CloudWatch~\cite{aws:cloudwatch} (server-side)} --- per Lambda,
        five metrics at 1-minute granularity: Invocations (Sum), Duration
        (Min/Avg/Max), Concurrent Executions (Max), Errors (Sum), and Throttles
        (Sum). Captures only what happens inside the Lambda execution
        environment, with no visibility into upstream queuing.
\end{itemize}

All timestamps were normalised to Europe/Rome (UTC+2, June~2026) before
cross-referencing the two datasets.

\subsection{Load generator placement}

Locust ran on a team member's personal laptop in Italy, connecting to the
three Lambda functions in \texttt{us-east-1} through their API Gateway HTTPS
endpoints over the public Internet. The reported response times therefore
include real-world network round-trip latency in addition to queueing and
Lambda execution time, which we consider representative of how a typical web
user would experience the service.

% ── Figure 1: server-side metrics, full width, top of page 2 ─────────────────
\begin{figure*}[t]
  \centering
  \begin{subfigure}[b]{0.48\linewidth}
    \centering
    \includegraphics[width=\linewidth]{concurrent_lambda_execution.png}
    \caption{Concurrent Executions over Time}
    \label{fig:cw_concurrent}
  \end{subfigure}
  \hfill
  \begin{subfigure}[b]{0.48\linewidth}
    \centering
    \includegraphics[width=\linewidth]{Duration_distribution_sidebyside.png}
    \caption{Duration Distribution per Operation}
    \label{fig:cw_dist}
  \end{subfigure}
  \caption{Server-side metrics from AWS CloudWatch.
    \textbf{(a)}~Concurrent Lambda executions over the full 3-hour session;
    peak concurrency of 9--11 is visible at the busiest scenarios and sits two
    orders of magnitude below the 1{,}000-concurrent production default,
    confirming that the Vocareum sandbox concurrency cap is the binding
    scalability constraint.
    \textbf{(b)}~Duration distribution per operation across all 72 scenarios;
    grayscale is the lightest ($\bar{x}$\,=\,143\,ms), edge detection the
    heaviest ($\bar{x}$\,=\,438\,ms), and the spread within each operation
    remains narrow, confirming that Lambda execution time is stable regardless
    of load level.}
  \label{fig:server_metrics}
\end{figure*}

% ══════════════════════════════════════════════════════════════════════════════
\section{Server-side Results (CloudWatch)}
\label{sec:server}

Figure~\ref{fig:server_metrics} and Table~\ref{tab:cloudwatch} present the
per-Lambda CloudWatch metrics aggregated over the full 3-hour test session.

\begin{table}[t]
\centering
\caption{Per-Lambda server-side summary (CloudWatch, full session).}
\label{tab:cloudwatch}
\small
\setlength{\tabcolsep}{3.5pt}
\begin{tabular}{lrrr}
\toprule
\textbf{Metric} & \textbf{resize} & \textbf{gray} & \textbf{edge} \\
\midrule
Total invocations       & 17{,}342 & 17{,}521 & 17{,}409 \\
Avg Duration (ms)       & 271.1    & 142.8    & 438.0    \\
p95 Duration (ms)       & 615.8    & 432.3    & 1{,}053.8 \\
Min Duration (ms)       & 11.2     & 1.6      & 8.5      \\
Max Duration (ms)       & 856.4    & 740.3    & 1{,}431.3 \\
Max concurrent exec.    & 11       & 10       & 9        \\
Total errors            & 0        & 0        & 0        \\
Min success rate (\%)   & 100      & 100      & 100      \\
Total throttles         & 0        & 0        & 0        \\
\bottomrule
\end{tabular}
\end{table}

Three observations stand out. First, Duration is \textbf{remarkably stable}
across the entire session: p95 sits at 432--1{,}054\,ms depending on the
operation and never spikes in response to increasing load, as visible in
Figure~\ref{fig:cw_dist}. This confirms that Lambda's execution environment
scales horizontally without degradation---each additional concurrent invocation
receives its own container and executes at the same speed as a lone request.
Second, \textbf{peak concurrent executions reached only 9--11}
(Figure~\ref{fig:cw_concurrent}), two orders of magnitude below the
1{,}000-concurrent production default~\cite{aws:concurrency}. The Vocareum
sandbox imposes an account-level concurrency cap that manifests as this ceiling;
its practical effect is that requests accumulate in the API Gateway queue rather
than spawning additional Lambda instances. This is the root cause of the
client-side failures discussed in Section~\ref{sec:client}. Third, \textbf{Lambda reported
zero errors, zero throttles, and 100\,\% success} across all 52{,}272
invocations. From the function's perspective, the service operated perfectly
throughout.

The Duration hierarchy is consistent with algorithmic complexity: grayscale
(143\,ms average) performs a single per-pixel luma conversion; resize (271\,ms)
applies a multi-tap LANCZOS resampling kernel; edge detection (438\,ms) runs two
sequential convolution passes. Because Lambda bills per millisecond of execution
time at the allocated memory, this ordering maps directly to a proportional cost
differential, quantified in Section~\ref{sec:cost}.

% ── Figure 2: client-side metrics, full width, bottom of page 2 ──────────────
\begin{figure*}[b]
  \centering
  \begin{subfigure}[b]{0.32\linewidth}
    \centering
    \includegraphics[width=\linewidth]{Client_p95_vs_concurrent users.png}
    \caption{Client p95 Latency vs.\ Users}
    \label{fig:locust_p95}
  \end{subfigure}
  \hfill
  \begin{subfigure}[b]{0.32\linewidth}
    \centering
    \includegraphics[width=\linewidth]{Throughput_vs_users.png}
    \caption{Throughput vs.\ Concurrent Users}
    \label{fig:locust_throughput}
  \end{subfigure}
  \hfill
  \begin{subfigure}[b]{0.32\linewidth}
    \centering
    \includegraphics[width=\linewidth]{Error_rate_vs_users.png}
    \caption{Client Error Rate (\%)}
    \label{fig:locust_errors}
  \end{subfigure}
  \caption{Client-side performance captured by Locust across all 72 scenarios.
    \textbf{(a)}~p95 response time grows quasi-linearly from $\approx$1\,s at
    1~user to $\approx$46\,s at 100~users for all three operations, despite
    Lambda execution remaining stable (cf.\ Table~\ref{tab:cloudwatch}).
    The three operation curves converge at high concurrency, confirming that
    the bottleneck is the upstream queue, not the function logic.
    \textbf{(b)}~Throughput saturates well below the nominal user count because
    each virtual user is blocked waiting for its previous response.
    \textbf{(c)}~HTTP failures are absent at 1 and 10~users, appear at 50, and
    reach 27\,\% for \textit{resize\_large} at 100~users; all originate in
    API Gateway, not in Lambda.}
  \label{fig:client_metrics}
\end{figure*}

% ══════════════════════════════════════════════════════════════════════════════
\section{Client-side Results (Locust)}
\label{sec:client}

Figure~\ref{fig:client_metrics} presents the three primary client-side metrics.
The contrast with the server-side picture is stark. While Lambda was completing
each individual invocation in under 1.5\,s (Table~\ref{tab:cloudwatch}),
end-to-end p95 response time grew quasi-linearly with the user count, reaching
$\approx$46\,s at 100~concurrent users for all three operations
(Figure~\ref{fig:locust_p95}). The three operation curves are nearly coincident
at high concurrency, indicating that operation type becomes irrelevant once the
API Gateway queue is saturated: a 143\,ms grayscale call and a 438\,ms edge
call both wait the same 45+ seconds in queue before they are even dispatched.

Throughput tells the same story from a supply-side perspective
(Figure~\ref{fig:locust_throughput}): completed requests per second plateau well
below the nominal user count because each Locust virtual user is serialised on
its own response. Adding virtual users beyond the saturation point increases
queue depth without increasing throughput, a classic symptom of a
queueing-limited system.

Across the 72~scenarios Locust recorded \textbf{315 failures}
(Figure~\ref{fig:locust_errors}): 161~HTTP~0 (TCP connection reset before any
response) and 154~HTTP~408 (request timeout returned by API Gateway after
29\,s). The worst single scenario was \textit{resize\_large at 100~users}, with
82~failures in the first repetition and 61 in the second (a 27\,\% failure
rate in the second run). Failures were entirely absent at 1 and 10~users and
first appeared at 50~users, consistent with the hypothesis that they are caused
by API Gateway queue exhaustion rather than any instability within Lambda.

The effect of payload size on client-perceived latency is summarised in
Table~\ref{tab:size}. At moderate load (10 users) the size effect is mild and
proportional to the payload; at peak load (100 users), larger images inflate
p95 response time by more than an order of magnitude across all three
operations, because the per-request service time amplifies into seconds of
additional queue wait for every other request behind it.

\begin{table}[h]
\centering
\caption{Client-side p95 response time (ms) by image size, at moderate
(10 users) and peak (100 users) load. S/M/L = small/medium/large images.
Values averaged across repetitions.}
\label{tab:size}
\small
\setlength{\tabcolsep}{3pt}
\begin{tabular}{l r r r r r r}
\toprule
& \multicolumn{3}{c}{\textbf{10 users}} & \multicolumn{3}{c}{\textbf{100 users}} \\
\cmidrule(lr){2-4} \cmidrule(lr){5-7}
\textbf{Op.} & \textbf{S} & \textbf{M} & \textbf{L} &
              \textbf{S} & \textbf{M} & \textbf{L} \\
\midrule
resize    & 320 & 2{,}000 & 12{,}000 & 3{,}050 & 41{,}000 & 96{,}500 \\
grayscale & 285 & 2{,}000 & 13{,}500 & 4{,}100 & 35{,}000 & 99{,}000 \\
edge      & 315 & 2{,}100 & 12{,}500 & 3{,}150 & 38{,}000 & 97{,}500 \\
\bottomrule
\end{tabular}
\end{table}

% ── Figure 3: cross-reference and cost, full width, top of page 4 ────────────
\begin{figure*}[t]
  \centering
  \begin{subfigure}[b]{0.50\linewidth}
    \centering
    \includegraphics[width=\linewidth]{End-to-end_vs_server-side_p95.png}
    \caption{End-to-End vs.\ Server-Side p95 Latency}
    \label{fig:xref_latency}
  \end{subfigure}
  \hfill
  \begin{subfigure}[b]{0.46\linewidth}
    \centering
    \includegraphics[width=\linewidth]{Cost_break-even_curve.png}
    \caption{6-Month Cost: Lambda vs.\ EC2~t3.small}
    \label{fig:cost_curve}
  \end{subfigure}
  \caption{\textbf{(a)}~End-to-end p95 (solid lines, Locust) vs.\ server-side
    p95 (dashed lines, CloudWatch) per operation and concurrency level.
    The gap between solid and dashed is the API Gateway queuing overhead; at
    100~users it accounts for $>$99\,\% of observed latency for all three
    operations. Note that the dashed lines for all three operations are nearly
    flat and indistinguishable near the x-axis, confirming that Lambda Duration
    does not increase with load.
    \textbf{(b)}~Six-month cost projection. Lambda cost grows linearly with
    invocation volume starting near zero; the always-on EC2~t3.small incurs a
    flat $\approx$\$546 regardless of traffic. Break-even occurs at
    $\approx$6--7\,M invocations per month, above which provisioned capacity
    becomes the rational economic choice.}
  \label{fig:analysis_economics}
\end{figure*}

% ══════════════════════════════════════════════════════════════════════════════
\section{Cross-reference: Where Does the Latency Live?}
\label{sec:cross}

Figure~\ref{fig:xref_latency} overlays the two p95 series---Locust end-to-end
(solid lines) and CloudWatch Lambda Duration (dashed lines)---for each
operation across all concurrency levels. The gap between solid and dashed
represents everything that happens outside pure Lambda execution: API Gateway
request routing, TCP connection establishment, payload serialisation, network
round-trip, and most importantly \textit{queue wait time} once the concurrency
cap is reached.

For each scenario we formalise this decomposition as:
\[
  \text{overhead} = \text{Locust\,p95} - \text{CloudWatch\,p95}
\]
which approximates the sum of all non-Lambda contributions to end-to-end
latency. Table~\ref{tab:xref} reports the result aggregated by operation and
user level, averaged across image sizes and repetitions.

\begin{table*}[h]
\centering
\caption{End-to-end latency decomposition (ms). Overhead = Locust\,p95 $-$
  CloudWatch\,p95. \%OH = overhead as a percentage of total end-to-end latency.
  Values averaged across the three image sizes and two repetitions per cell.}
\label{tab:xref}
\small
\begin{tabular}{llrrrr}
\toprule
\textbf{Operation} & \textbf{Users} &
\textbf{Locust p95 (ms)} & \textbf{CW p95 (ms)} &
\textbf{Overhead (ms)} & \textbf{\%OH} \\
\midrule
resize    &   1 &  1{,}097 &  272 &    826 & 75\,\% \\
resize    &  10 &  4{,}773 &  274 &  4{,}499 & 94\,\% \\
resize    &  50 & 28{,}133 &  272 & 27{,}861 & 99\,\% \\
resize    & 100 & 46{,}850 &  270 & 46{,}580 & 99.4\,\% \\
\midrule
grayscale &   1 &  1{,}323 &  189 &  1{,}134 & 86\,\% \\
grayscale &  10 &  5{,}262 &  178 &  5{,}083 & 97\,\% \\
grayscale &  50 & 31{,}492 &  176 & 31{,}316 & 99\,\% \\
grayscale & 100 & 46{,}033 &  175 & 45{,}858 & 99.6\,\% \\
\midrule
edge      &   1 &  1{,}660 &  439 &  1{,}221 & 74\,\% \\
edge      &  10 &  4{,}972 &  446 &  4{,}525 & 91\,\% \\
edge      &  50 & 27{,}618 &  441 & 27{,}177 & 98\,\% \\
edge      & 100 & 46{,}216 &  441 & 45{,}775 & 99.0\,\% \\
\bottomrule
\end{tabular}
\end{table*}

Even at a single concurrent user, 74--86\,\% of end-to-end latency is overhead
unrelated to Lambda execution, reflecting fixed costs of TCP handshake, TLS
negotiation, and API Gateway routing. By 10~users the overhead fraction rises to
91--97\,\%; by 50--100~users it exceeds 99\,\% for every operation. At this
point Lambda Duration---though perfectly stable---is effectively invisible to
the end user. One practical implication: optimising the Lambda
function code (e.g., switching LANCZOS to a faster interpolation filter) would
produce \textit{no measurable improvement} in user-facing latency under any
non-trivial load, because the bottleneck is the upstream queue, not the function
itself. The correct remedy is to raise the account-level concurrency limit or
adopt AWS Lambda Provisioned Concurrency to pre-warm containers.

This behaviour is consistent with Little's Law applied to the upstream queue.
With an arrival rate $\lambda$ that grows with the number of concurrent users
and a service capacity bounded by the observed concurrency ceiling
$C \approx 10$, the queue settles to a steady-state length of roughly
$\lambda \cdot \overline{W}$, where $\overline{W}$ is the per-request service
time. Once $\lambda$ exceeds $C / \overline{W}$ the queue saturates and the
wait time $W = N / \lambda$ grows linearly with the backlog~$N$, eventually
triggering the API Gateway 29\,s integration timeout. This is precisely the
regime in which the HTTP~408 failures of Section~\ref{sec:client} appear.

\subsection{Two perspectives on success}

The coexistence of a perfect CloudWatch record with 315 client-side failures is
not a contradiction. It reflects that the two systems count fundamentally
different events. CloudWatch counts Lambda invocations that started and completed
within the function timeout. Locust counts every HTTP response that was not a
2xx, regardless of which layer produced it. All 315~failures occurred upstream
of Lambda: 161~HTTP~0 responses indicate API Gateway closed the TCP connection
without forwarding the request to any Lambda (queue exhausted, connection
reset); 154~HTTP~408 responses indicate the 29-second integration
timeout~\cite{aws:apigw-limits} fired while the request was still waiting in
queue. The longest individual Lambda execution was 1{,}431\,ms for edge
detection. The 30-second function timeout was never approached across any of
the 52{,}272 invocations. A deployment that exposes only Lambda metrics would
therefore present a falsely healthy picture of a system that, at peak load, was
rejecting more than one in four user requests at the gateway.

% ══════════════════════════════════════════════════════════════════════════════
\section{Cost Projection --- Lambda vs.\ EC2}
\label{sec:cost}

Figure~\ref{fig:cost_curve} visualises the six-month cost comparison between
the serverless deployment and an equivalent always-on EC2 baseline. Using the
measured average Durations (resize 271\,ms, grayscale 143\,ms, edge 438\,ms;
mean 284\,ms across operations) we apply public on-demand
pricing~\cite{aws:lambda-pricing,aws:ec2-pricing}:

\begin{itemize}
  \item \textbf{Lambda}: \$0.20 per million invocations plus
        \$1.67\,$\times$\,10$^{-5}$ per GB-second at 512\,MB allocation.
        The first 1\,M invocations per month and 400{,}000\,GB-seconds are
        covered by the AWS Free Tier indefinitely.
  \item \textbf{EC2~t3.small}: \$0.0208\,/\,hour on-demand in
        \texttt{us-east-1} $\approx$~\$15.18/month, regardless of request volume,
        running 24/7. Over six months this totals $\approx$~\$91. AWS Free Tier
        covers 1\,M Lambda invocations and 400\,000\,GB-seconds monthly.
\end{itemize}

Table~\ref{tab:cost} materialises the calculation at representative monthly
volumes, applying the formula above to the mean measured Duration of 284\,ms
and 512\,MB memory.

\begin{table}[h]
\centering
\caption{Six-month deployment cost at representative monthly volumes.
Lambda cost computed as
$(\$0.20\cdot N_{\text{bill}}/10^6 + \$1.67\!\times\!10^{-5}\cdot GB{\text{-}}s_{\text{bill}})$,
where $N_{\text{bill}}$ = billable invocations (after free tier), $GB{\text{-}}s_{\text{bill}}$ = billable GB-seconds (after free tier).}
\label{tab:cost}
\small
\setlength{\tabcolsep}{4pt}
\begin{tabular}{r r r l}
\toprule
\textbf{Inv./month} & \textbf{Lambda (6m)} & \textbf{EC2 (6m)} & \textbf{Winner} \\
\midrule
   100\,K & \$0.0     & \$91 & Lambda \\
     1\,M & \$0.0     & \$91 & Lambda \\
     5\,M & \$0.0     & \$91 & Lambda \\
     7\,M & \$91      & \$91 & break-even \\
    10\,M & \$146     & \$91 & EC2 \\
   100\,M & \$1{,}887 & \$91 & EC2 \\
\bottomrule
\end{tabular}
\end{table}

Lambda cost starts near zero and grows linearly with traffic; EC2 cost is a
flat step function. The \textbf{break-even} falls at roughly
\textbf{4.7\,M invocations per month}: below that threshold Lambda is
economicly rational; above it, a provisioned instance becomes cheaper, provided
a single t3.small can sustain the required response-time SLA at the
corresponding request rate, a condition the present study has not verified. For
mixed-operation workloads the break-even shifts depending on the fraction of
edge-detection calls, which cost 3$\times$ more in Lambda Duration than
grayscale.

% ══════════════════════════════════════════════════════════════════════════════
\section{Discussion and Limitations}
\label{sec:discussion}

\textbf{Vocareum sandbox concurrency cap.}
The account-level concurrency ceiling clamped peak concurrent executions to
9--11, two orders of magnitude below the 1{,}000-concurrent production
default~\cite{aws:concurrency}. This is the direct cause of the API Gateway
queue saturation observed from 50~users onward. In a production account with
the default limit, the same 100-user load would spawn up to 100 parallel Lambda
containers, each completing in under 1.5\,s, with no queue accumulation and no
client-visible failures. The sandbox thus stress-tests the upstream queuing
layer rather than Lambda's own scalability, and the latency and failure numbers
cannot be directly extrapolated to a real-world deployment.

\textbf{Load generator over the public Internet.}
Locust ran on a personal laptop in Italy reaching the Lambda endpoints in
\texttt{us-east-1} over the public Internet, so the reported response times
already include the round-trip network latency that a real user would also
experience. Our overhead values (Table~\ref{tab:xref}) therefore embed a
fixed RTT component that is independent of Lambda or queueing. Users with
different geographic distance to \texttt{us-east-1} would see proportionally
different base latency.

\textbf{Operation cost proportional to CPU complexity.}
Edge detection ran 62\,\% longer than resize and 207\,\% longer than grayscale
on average. Because Lambda bills per millisecond at the allocated memory, this
translates directly into proportional cost differentials. A production deployment
processing a mixed workload should weight the cost projection by the expected
operation mix rather than using a single average Duration.

\textbf{Missing cold-start measurement.}
A dedicated cold-start probe was excluded due to budget constraints. Such a
probe would invoke each Lambda after a $\geq$20-minute idle period to force
cold container initialisation, then record Init Duration from CloudWatch Logs
Insights. Init Duration
is explicitly cited in the course brief and its absence is a recognised gap.
Future work should reserve $\approx$5\,\% of the credit budget for 10~cold-start
observations per Lambda, sufficient to characterise initialisation overhead and
its contribution to worst-case p99 latency.

% ══════════════════════════════════════════════════════════════════════════════
\section{Conclusions}
\label{sec:conclusions}

This study deployed and evaluated three image-processing Lambda functions under
a controlled factorial load-test. The four principal findings are:

\begin{enumerate}
  \item \textbf{Server-side health does not imply user-side health.}
        Lambda reported zero errors and 100\,\% success across 52{,}272
        invocations, yet the same workload produced 46\,s p95 latency and
        27\,\% HTTP failures for end users at 100~concurrent users. All failures
        originated in API Gateway and are invisible to Lambda-only monitoring.
  \item \textbf{Queueing overhead dominates at scale.}
        At 50--100 concurrent users, API Gateway queueing accounts for
        $>$99\,\% of observed end-to-end latency (Table~\ref{tab:xref},
        Figure~\ref{fig:xref_latency}). Code-level Lambda optimisations have
        negligible impact; the correct remedy is raising the concurrency limit
        or adopting Provisioned Concurrency.
  \item \textbf{Operation type determines cost, not user-facing latency at scale.}
        Grayscale, resize, and edge differ by up to 3$\times$ in Lambda
        Duration---and therefore billing cost---but their client-side p95 values
        converge to $\approx$46\,s at high concurrency because the bottleneck
        shifts entirely to the upstream queue (Figure~\ref{fig:locust_p95}).
  \item \textbf{Lambda is cost-effective below $\approx$7\,M invocations/month.}
        Above that threshold an always-on EC2~t3.small is cheaper, subject to
        latency SLA verification at the corresponding load
        (Figure~\ref{fig:cost_curve}).
\end{enumerate}

The methodological lesson is that evaluating a serverless deployment
requires \textit{both} a server-side and a client-side perspective. A Locust
load generator combined with CloudWatch metric export, cross-referenced as shown
in Section~\ref{sec:cross}, is the minimum instrumentation needed to arrive at
a correct and complete characterisation of system behaviour.

% ══════════════════════════════════════════════════════════════════════════════
\section*{References}

\begingroup
\small
\renewcommand{\section}[2]{}
\begin{thebibliography}{9}

\bibitem{aws:lambda-pricing}
Amazon Web Services.
\textit{AWS Lambda Pricing}.
\url{https://aws.amazon.com/lambda/pricing/}
Accessed June 2026.

\bibitem{aws:concurrency}
Amazon Web Services.
\textit{Managing AWS Lambda Function Concurrency}.
AWS Lambda Developer Guide.
\url{https://docs.aws.amazon.com/lambda/latest/dg/lambda-concurrency.html}

\bibitem{aws:apigw-limits}
Amazon Web Services.
\textit{Amazon API Gateway Quotas and Important Notes}.
\url{https://docs.aws.amazon.com/apigateway/latest/developerguide/limits.html}

\bibitem{aws:cloudwatch}
Amazon Web Services.
\textit{Using Amazon CloudWatch Metrics with AWS Lambda}.
\url{https://docs.aws.amazon.com/lambda/latest/dg/monitoring-metrics.html}

\bibitem{aws:ec2-pricing}
Amazon Web Services.
\textit{Amazon EC2 On-Demand Pricing --- t3.small, us-east-1}.
\url{https://aws.amazon.com/ec2/pricing/on-demand/}

\bibitem{pillow:docs}
Pillow Contributors.
\textit{Pillow (PIL Fork) Documentation --- Image Filters and Resize}.
\url{https://pillow.readthedocs.io/}
Accessed June 2026.

\bibitem{locust:docs}
Locust Contributors.
\textit{Locust --- A Modern Load Testing Framework}.
\url{https://locust.io/}
Accessed June 2026.

\end{thebibliography}
\endgroup

\end{document}
